% !TeX program = xelatex
\documentclass[t]{beamer}
\usepackage{graphicx}
\usepackage{amsthm}
\usepackage{amsmath}
\usepackage{amssymb}
\usepackage{tikz,pgfplots}
\pgfplotsset{compat=1.18}
\usetikzlibrary{positioning, automata}
\usepackage{hyperref}
\usepackage{algorithm}
\usepackage{algpseudocode}

% Remove all \End commands
\algnotext{EndIf}
\algnotext{EndFor}
\algnotext{EndWhile}
\algnotext{EndLoop}
\algnotext{EndForAll}
\algnotext{EndRepeat}
\algnotext{EndProcedure}
\algnotext{EndFunction}
\algnotext{Function}

\usepackage{tcolorbox}
\usepackage{xcolor}
\usepackage{fontspec}
\usepackage{libertinus}
\usepackage[greek, english]{babel}
\usepackage{mathtools}
\usetikzlibrary{arrows,positioning, calc}
\tikzstyle{vertex}=[draw,fill=white,circle,minimum size=20pt,inner sep=0pt]
\usetikzlibrary{shapes.geometric, arrows}
\usepackage{listings}

% Beamer presentation setup
\usetheme{Madrid}
\usecolortheme{default}
\setbeamertemplate{navigation symbols}{}

\usetikzlibrary{shapes,decorations,decorations.pathreplacing,arrows,calc,arrows.meta,fit,positioning}

\tikzset{
	auto,node distance =1 cm and 1 cm, semithick,
	state/.style ={ellipse, draw, minimum width = 0.7 cm},
	point/.style = {circle, draw, inner sep=0.04cm,fill,node contents={}},
	bidirected/.style={Latex-Latex,dashed},
	directed/.style={-Latex},
	el/.style = {inner sep=2pt, align=left, sloped}
}

\let\MakeUppercase\relax

% Redefine Function command to omit the "Function" prefix - start only (no end)
\algdef{S}[FUNCTION]{Function}%
   [2]{\textproc{#1}\ifthenelse{\equal{#2}{}}{}{(#2)}}%

% Redefine line number presentation
\algrenewcommand{\alglinenumber}[1]{\ifnum#1=0\else\footnotesize #1\fi}

% Define custom colors
\definecolor{myyellow}{RGB}{255,255,150}

% Define algorithmic environment inside a yellow-colored box
\newenvironment{yellowalgo}{%
	\begin{tcolorbox}[colback=myyellow, colframe=myyellow, boxrule=0.5pt]
	\begin{algorithmic}[1]
	\setcounter{ALG@line}{-1}
}{%
	\end{algorithmic}
	\end{tcolorbox}
}

\setbeamercovered{transparent} 

% Custom command for function names
\newcommand{\Func}[1]{\textproc{#1}}

\renewcommand{\L}{\mathcal{L}}
\newcommand{\Create}{\textbf{Create} }
\DeclareMathOperator{\argmax}{argmax}
\DeclareMathOperator{\argmin}{argmin}
\newcommand{\True}{\texttt{True}}
\newcommand{\False}{\texttt{False}}

% Redefine the theorem environment to use normal font for the body text
\makeatletter
\def\th@plain{%
	\thm@notefont{}%
	\normalfont
}
\makeatother

\newtheorem{fig}{Figure}
\newtheorem{obs}{Παρατήρηση}
\newtheorem{cor}{Πόρισμα}

\title{Προβλήματα στα Συντομότερα Μονοπάτια}
\subtitle{Θεωρία Γραφημάτων και Εφαρμογές}
\author{}
\date{}

\author[Κολιόπουλος, Βογιατζής]{Α. ~Κολιόπουλος\and Δ.~Βογιατζής}

\institute[ΠΑΔΑ]{Πανεπιστήμιο Δυτικής Αττικής}

\date[Θεωρία Γραφημάτων και Εφαρμογές]%{Παρουσίαση Εφαρμογών}


\AtBeginSection[]
{
  \begin{frame}
    \frametitle{Προβλήματα}
    \tableofcontents[currentsection]
  \end{frame}
}

\begin{document}

\setcounter{framenumber}{-1}


\begin{frame}[plain]
\titlepage
\vfill
\begin{center}
	\includegraphics[width=0.28\textwidth]{UniWA logo.png}
\end{center}
\end{frame}


% \begin{frame}
% \frametitle{Προβλήματα}
% \tableofcontents
% \end{frame}

\section{Σημαντικοί Δρόμοι - $28$ος ΠΔΠ Καμπ (seniors)}

\begin{frame}
    \frametitle{Σημαντικοί Δρόμοι - $28$ος ΠΔΠ Καμπ (seniors)}
    \underline{Είσοδος:}
	\pause
	\vspace{0.25cm} 
	\begin{itemize}
		\item Κατευθυνόμενο γράφημα $G = \left(V.E\right) : \forall e \in G.E, \; w(e) \in \mathbb{N}_{+}$. 
	\end{itemize}
	\pause
	\vspace{0.25cm}
	\underline{Ιδιότητες:}
	\pause
	\vspace{0.25cm}
	\begin{itemize}
		\item Η κορυφή με τον αριθμό $1$ αποτελεί την αφετηρία.
		\pause
		\item Οποιαδήποτε άλλη κορυφή είναι προσπελάσιμη από την αφετηρία. \label{prop2}
		\pause
		\item \emph{Σημαντική} είναι κάθε ακμή που η μείωση του βάρους της κατά $1$, επιφέρει αντίστοιχη μείωση στην απόσταση της αφετηρίας προς κάποια άλλη κορυφή.
		\pause
		\item \emph{Πολύ σημαντική} είναι κάθε ακμή που η αύξηση του βάρους της κατά $1$, επιφέρει αντίστοιχη αύξηση στην απόσταση της αφετηρίας προς κάποια άλλη κορυφή.
	\end{itemize}
\end{frame}

\begin{frame}
	\frametitle{Σημαντικοί Δρόμοι - $28$ος ΠΔΠ Καμπ (seniors)}
	\underline{Έξοδος:}
	\pause
	\vspace{0.25cm} 
	\begin{itemize}
		\item $\#$Σημαντικών ακμών.
		\pause 
		\item $\#$Πολύ σημαντικών ακμών.
	\end{itemize}
	\pause
	\vspace{0.25cm} 
	\underline{Λύση:}
	\pause
	\vspace{0.25cm} 
	\begin{itemize}
		\item Εκτελούμε τον αλγόριθμο του Dijkstra και παράγουμε το ΚΑΓ συντομότερων μονοπατιών, έστω $\hat{G}$.
		\pause
		\item Προσθέτουμε κάθε ακμή του $\hat{G}$ στο σύνολο των σημαντικών ακμών, ονομαστικά $I$.
		\pause
		\item Προσθέτουμε κάθε ακμή $(u, v) \in \hat{G}$, για την οποία ισχύει ότι $|\hat{G}.N^-(v) = 1|$, στο σύνολο των πολύ σημαντικών ακμών, έστω $I'$.
		\pause
		\item Εκτυπώνουμε τα μεγέθη των $I$, $I'$.
	\end{itemize}
\end{frame}


% \begin{frame}
% 	\frametitle{Σημαντικοί Δρόμοι - $28$ος ΠΔΠ Καμπ (seniors)}
	
% 	\begin{yellowalgo}
% 		\Function{Critical}{$G$}
%     		\For{\textbf{all} $v \in G.V \setminus \{1\}$}
%         		\State $D[v] \gets \infty;$
%     		\EndFor
%     		\State $D[1] \gets 0; \; S \gets \emptyset;$
%     		\While{$S \neq G.V$} 
%         		\State $v \gets {\arg\min}\left\{D[u] : u \in G.V \setminus S\right\}; \; S \gets S \cup \{v\};$
%        			\For{\textbf{all} $u \in G.N^+(v) \setminus S$}
%             		\If{$D[v] + w(v, u) < D[u]$}
%                 		\State $D[u] \gets D[v] + w(v, u); \; \hat{G}.N^-(u) \gets \{v\};$
% 					\ElsIf{$D[v] + w(v, u) = D[u]$}
% 						\State $\hat{G}.N^-(u) \gets \hat{G}.N^-(u) \cup \{v\};$
%             		\EndIf
%         		\EndFor
%     		\EndWhile
% 			\For{$v \in G.V$}
% 			\EndFor
% 		\EndFunction
% 	\end{yellowalgo}
	
% \end{frame}

\begin{frame}
	\frametitle{Σημαντικοί Δρόμοι - $28$ος ΠΔΠ Καμπ (seniors)}
	\underline{Ορθότητα:}
	\pause
	\begin{block}{Σημαντικές ακμές}
		Οι σημαντικές ακμές είναι αυτές του ΚΑΓ που παράγει ο Dijkstra. Ισοδύναμα, $I = \hat{G}.E$.
	\end{block}
	\pause
	\begin{block}{\emph{Απόδειξη.}}
		Έστω ότι $I \neq \hat{G}.E$. Διακρίνουμε τις εξής δύο περιπτώσεις. 
		\pause
		\\
		\vspace{0.25cm} 
		(Πρώτη περίπτωση). Έστω ότι υπάρχει ακμή $(v, u) \in G.E$, εκτός ΚΑΓ, η οποία είναι σημαντική. 
		Γνωρίζουμε σίγουρα ότι $u \neq 1$, αφού $d(1) = 0$ και $w(e) \geq 1$, για κάθε $e \in G.E$.
		Δεδομένου ότι κάθε κορυφή είναι προσπελάσιμη από την αφετηρία (Ιδιότητα \ref{prop2}), έπεται ότι υπάρχει 
		συντομότερο μονοπάτι από την κορυφή $1$ στη $u$. Έστω $v^*$ η τελευταία κορυφή σε ένα από τα συντομότερα μονοπάτια 
		που καταλήγουν στην $u$.
	\end{block}
\end{frame}

\begin{frame}
	\frametitle{Σημαντικοί Δρόμοι - $28$ος ΠΔΠ Καμπ (seniors)}
	\begin{block}{\emph{Απόδειξη. (συνέχεια)}}
		 Έχουμε ότι 
		 \pause
		\begin{align*}
			d(v) + w(v, u) > d(v^*) + w(v^*, u) &\Rightarrow d(v) + w(v, u) - 1 \geq d(v^*) + w(v^*, u) \\ %=d(u) \\
			&\Rightarrow d(v) + w(v, u) - 1 \geq d(u).
		\end{align*}
		\pause
		Επομένως δεν υπάρχει ακμή εκτός του $\hat{G}$, η οποία είναι σημαντική.
		\pause
		\\
		\vspace{0.25cm} 
		(Δεύτερη περίπτωση). Έστω ότι υπάρχει ακμή $(v, u) \in \hat{G}.E$, που δεν είναι σημαντική. %Δεδομένου ότι η ίδια ανήκει στο $\hat{G}.E$, γνωρίζουμε 
		%ότι υπάρχει σε τουλάχιστον ένα συντομότερο μονοπάτι από 
	\end{block}
\end{frame}

% \begin{frame}
% \begin{center}
% \scalebox{1.5}{
% \begin{tikzpicture}
%     \node[state] (1) at (0,0) {$1$};
%     \node[state] (2) at (1.5, -2) {$2$};
%     \node[state] (3) at (-1.5, -2) {$3$}; %[below=of 2] {$3$}
%     % \draw[directed] (1) to[bend left = 20] node[above] {$3$}(2);
%     % \draw[directed] (2) to[bend left = 20] node[below] {$5$}(1);
%     % \draw[directed] (1) -- node[left, xshift=-2pt] {$7$} (3);
%     % \draw[directed] (2) -- node[right, xshift=2pt] {$1$} (3);
% \end{tikzpicture}
% }
% \end{center}
% \end{frame}

\section{Σύστημα πλοήγησης}

\end{document}




	% \begin{yellowalgo}
	% 	\Function{Critical}{$G$}
	% 		\State $\hat{G} \gets \Func{Dijkstra}(G, 1);$
	% 		\State $I \gets \hat{G}.E;$ %\Comment{Σημαντικοί δρόμοι.}
	% 		\State $I' \gets \left\{(u, v) \in \hat{G}.E : \left|N^-(v)\right| = 1\right\}$ %\Comment{Πολύ σημαντικοί δρόμοι.}

	% 		\State \textbf{print} $|I|, \left|I'\right|$ 
	% 	\EndFunction
	% \end{yellowalgo}

% \underline{Αλγόριθμος:}
